\documentclass[10pt,letterpaper]{article}
\usepackage[margin=0.50in]{geometry}
\usepackage[T1]{fontenc}
\usepackage[scaled=0.94]{tgheros}
\usepackage[hidelinks]{hyperref}
\usepackage{enumitem}
\usepackage{titlesec}
\usepackage{xcolor}
\usepackage{tabularx}
\usepackage{array}

\definecolor{ink}{HTML}{171717}
\definecolor{muted}{HTML}{555555}
\definecolor{rule}{HTML}{B8B8B0}

\pagestyle{empty}
\setlength{\parindent}{0pt}
\setlength{\parskip}{0pt}
\setlist[itemize]{leftmargin=1.2em,itemsep=2.5pt,topsep=2pt,parsep=0pt,partopsep=0pt}
\newcommand{\headingfont}{\sffamily}
\titleformat{\section}{\headingfont\Large\bfseries\color{ink}}{}{0pt}{}[\vspace{2pt}\color{rule}\titlerule]
\titlespacing*{\section}{0pt}{10pt}{6pt}
\renewcommand{\familydefault}{\sfdefault}

\newcommand{\resumeRole}[4]{
  \begin{tabular*}{\textwidth}{@{\extracolsep{\fill}} l r}
    \textbf{#1} & #2 \\
    \textit{#3} & \textit{#4}
  \end{tabular*}
}

\newcommand{\educationRow}[4]{
  \begin{tabularx}{\textwidth}{@{}>{\bfseries}X l l r@{}}
    #1 & #2 & GPA: #3 & #4
  \end{tabularx}
}

\begin{document}

\begin{center}
  {\Huge \textbf{Yuchen Liu}}\\[-1pt]
  Seattle, WA \quad | \quad
  \href{mailto:yql6113@gmail.com}{yql6113@gmail.com} \quad | \quad
  814-876-2346 \quad | \quad
  \href{https://github.com/Upcccccc}{github.com/Upcccccc} \quad | \quad
  \href{https://www.linkedin.com/in/liuyuche/}{linkedin.com/in/liuyuche}\\
  \vspace{2pt}
  \textcolor{muted}{AI systems engineer building RL training infrastructure, search/data pipelines, and research-grade open-source tools.}
  \vspace{8pt}
\end{center}

\section*{Education}
\educationRow{M.S. Computer Science}{University of Pennsylvania, Philadelphia}{3.80/4.00}{May 2025}
\educationRow{B.S. Applied Economics}{Pennsylvania State University, State College}{3.86/4.00}{May 2023}

\section*{Experience}
\resumeRole{Software Engineer, Search Infrastructure}{Aug 2025 -- Present}{DoorDash}{Seattle, WA}
\begin{itemize}
  \item Built \textbf{Notus}, a next-generation search ingestion platform using \textbf{Iceberg, Kafka, and Spark} to decouple document ingestion from serving, enable configuration-driven indexing, and unblock schema evolution for hybrid search.
  \item Built \textbf{Vela}, Argo's hybrid retrieval platform, enabling \textbf{Lucene-native hybrid search} and federated retrieval with vector backends such as \textbf{Milvus} behind a single API.
  \item Migrated \textbf{30+ production stacks} to modular \textbf{Jsonnet}, eliminating roughly \textbf{56K lines} of duplicate configuration and reducing config verification time by \textbf{90\%}.
  \item Stabilized production indexers with \textbf{500GB EBS storage}, pre-deployment validators, and \textbf{Kyverno guardrails} after RCA for \textbf{\$21K revenue-risk incidents}, preventing capacity-related outages.
\end{itemize}
\vspace{4pt}

\resumeRole{Software Engineering Intern}{Oct 2024 -- May 2025}{Penn Medicine TissueLab}{Philadelphia, PA}
\begin{itemize}
  \item Developed a \textbf{Python/FastAPI} microservice platform to orchestrate \textbf{PyTorch} models for medical image analysis, enabling natural-language-driven segmentation and classification workflows.
  \item Implemented an event-driven \textbf{DAG workflow engine} for async task execution, processing \textbf{10K+ pathology images daily} and streaming \textbf{12M+ imaging events} through WebSockets with sub-10ms latency.
  \item Integrated generative and discriminative deep learning models into a unified service layer with \textbf{asynchronous APIs} for cross-department usage.
\end{itemize}
\vspace{4pt}

\resumeRole{Software Engineering Intern}{Apr 2024 -- Aug 2024}{Information Technology of CAS}{Remote}
\begin{itemize}
  \item Designed \textbf{Kafka and Redis-backed} event streaming services that reduced \textbf{p95 API latency from 2s to 200ms} while processing millions of IoT events.
  \item Contributed to distributed \textbf{Spring Boot} monitoring services, delivered \textbf{35+ REST APIs}, and maintained \textbf{99.7\% uptime} in Unix-based production environments.
\end{itemize}

\section*{Open Source}
\href{https://github.com/RiddleHe/nanochat}{\textbf{nanoRL}}
\begin{itemize}
  \item Led an open-source \textbf{RL training framework} inside nanochat for clean objective definitions, rollout experiments, RL algorithm research, reproducible ablations, and \textbf{vLLM inference serving}.
\end{itemize}
\vspace{2pt}
\href{https://github.com/RiddleHe/nanochat}{\textbf{nanochat}}
\begin{itemize}
  \item Built a hackable pretraining and chat stack supporting architecture definition and \textbf{FLOP-controlled model ablations}.
\end{itemize}
\vspace{2pt}
\href{https://github.com/RiddleHe/llm-interp}{\textbf{llm-interp}}
\begin{itemize}
  \item Authored reproducible interpretability scripts for model circuit research, including \textbf{attention-sink analysis} and reproductions of \textbf{LLM decode indeterminism} findings.
\end{itemize}

\section*{Research}
\href{https://arxiv.org/abs/2606.02780}{\textbf{Do Value Vectors in Deep Layers Need Context from the Residual Stream?}} \hfill EMNLP 2026, under review\\
\begin{itemize}
  \item Co-first author. Challenged the standard attention mechanism of computing value vectors from the residual stream, discovering that deep layers benefit from \textbf{context-free value vectors} that preserve original token information.
  \item Proposed \textbf{Bank of Values}, a learned value-vector table for the last third of layers that eliminates the \textbf{V cache} and improves validation loss and average scores across \textbf{21 benchmarks} on \textbf{135M and 780M} models.
\end{itemize}
\vspace{2pt}
\href{https://arxiv.org/abs/2509.21576}{\textbf{Vision Language Models Cannot Plan, but Can They Formalize?}} \hfill ECCV 2026, under review\\
\begin{itemize}
  \item Studied whether \textbf{VLMs} can formalize visual planning tasks into solver-executable \textbf{PDDL problem files}.
\end{itemize}
\vspace{2pt}
\href{https://arxiv.org/abs/2509.20279}{\textbf{A co-evolving agentic AI system for medical imaging analysis}} \hfill Nature, under review\\
\begin{itemize}
  \item Built agentic medical-imaging workflows for \textbf{tool selection, planning, knowledge updates}, and human-in-the-loop co-evolution.
\end{itemize}

\section*{Skills}
\textbf{Languages:} Python, Kotlin, Java, C++, JavaScript;
\textbf{ML/Infra:} PyTorch, vLLM, Spark, Kafka, Iceberg, Delta Lake, FastAPI, Redis, Docker, Kubernetes, AWS, Spring Boot

\end{document}
